An essential process in conceptual hydrological modeling is determining the optimal model parameter, corresponding to the global maximum, through calibration and validation. Depending on the model's complexity, the number of required parameters varies widely, and so does the intricacy of finding the optimum. Multiple equally valuable optimized solutions are also within the realm of possibility \autocite{Beven.2012}. Several methodologies are available for parameter optimization, for example, the Markov Chain Monte Carlo (MCMC)  Method and the Generalized Likelihood Uncertainty Estimation (GLUE) method. While not all techniques are suitable for every model, a range of options exists to address different scenarios \autocite{SPOTting}. This analysis focuses on the \ac{de} scheme. \\ 

Given the constraints of the project timeline, the focus will be only on calibration processes, and the results will serve as a foundation for the following assignments. The process of defining the parameter ranges for the solute search will not be addressed, but it is acknowledged as a complex challenge. If the parameter range is chosen incorrectly, the optimized solution may be excluded before the optimization process begins. Additionally, certain model processes will be deactivated after finding the optimal \ac{mp}. This provides a greater understanding of the model and input data. \\  


\begin{comment}
• Global model parameter optimization using Differential Evolution
• High flow event
• Get a feeling of how model parameter optimization works
• Various objective functions
• Reference: Long-term lumped projections of groundwater balances in the face of limited
data, Ejaz 2024, Ph.D. Thesis
\end{comment}